\section{بازتعریف مسئله و مبانی نظری}

\begin{frame}{مسئله‌ی روزنامه‌فروش}
  \begin{itemize}
    \item ۲۰۰ نفر رزرو کرده‌اند. اگر ۲۰۰ پرس بپزی و ۱۵ نفر نیایند، ۱۵ پرس دور ریخته
      می‌شود، حدود $5.6$ میلیون تومان
    \item اگر ۱۸۵ پرس بپزی و ۱۹۰ نفر بیایند، ۵ دانشجو غذا ندارند — کدام بدتر است؟
    \item جواب کلاسیک (Newsvendor): $\displaystyle \tau^{*} = \frac{C_u}{C_u+C_o}$
      — مقدار بهینه میانگین تقاضا نیست، چندک بحرانی توزیع تقاضاست
  \end{itemize}
  \begin{figure}
    \centering
    \includegraphics[width=0.72\textwidth]{slides_03_asymmetric_cost}
  \end{figure}
  \note{این همان مسئله‌ای است که در پژوهش در عملیات با نام Newsvendor شناخته می‌شود.
  جمله‌ی کلیدی: اگر کمبود چهار برابر مازاد گران است، آن‌قدر بپز که فقط در $20\%$ مواقع کم
  بیاوری.}
\end{frame}

\begin{frame}{شکست فرمول‌های کلاسیک: نوسان واقعی تا $15.6$ برابر فرضیات نظری است}
  \begin{itemize}
    \item \textbf{رد فرض استقلال افراد:} فرمول‌های تحلیلی غیبت را مستقل می‌دانند (توزیع دوجمله‌ای)، اما رفتار جمعی نوسان را تا $15.6$ برابر فراتر از تئوری می‌برد
    \item \textbf{ماندگاری نوسان پس از کنترل ساختاری:} حتی با کنترل سلف، وعده، تقویم و روز هفته، همچنان $7.6$ برابر نوسان اضافه باقی می‌ماند
    \item \textbf{پیامد راهبردی:} در میان ۷٬۵۷۹ سلول ناهمگن، چندک بهینه فرمول بسته ندارد و باید \textbf{مستقیماً از داده آموخته شود}
  \end{itemize}
  \begin{figure}
    \centering
    \includegraphics[width=0.55\textwidth]{4.8_variance_vs_res}
  \end{figure}
  \note{در رویکرد کلاسیک پژوهش عملیات، ابتدا یک توزیع پارامتری (مثل دوجمله‌ای با فرض استقلال
  افراد) برازش می‌شود و سپس چندک بهینه تحلیلی حساب می‌شود. اما داده‌ها نشان دادند فرض استقلال
  کاملاً باطل است و نوسان واقعی بین $3.7$ تا $15.6$ برابر فرمول نظری است (بیش‌پراکندگی صعودی).
  حتی کنترل کامل زمان و مکان هم این شکاف را پر نمی‌کند ($7.6$ برابر نوسان اضافه باقی می‌ماند).
  ضمناً ۷٬۵۷۹ سلول مستقل با شرایط متفاوت داریم. نتیجه: چندک نباید از فرمول بسته بیرون بیاید،
  باید مستقیماً از داده یاد گرفته شود. شاهد: \lr{F07}، \lr{F08}، \lr{F09}.}
\end{frame}

\begin{frame}{نیوزوندور داده‌محور و رگرسیون کوانتایل}
  \begin{itemize}
    \item به‌جای «اول توزیع را تخمین بزن، بعد بهینه کن»، مستقیماً تصمیم از داده یاد گرفته
      می‌شود (Ban \& Rudin 2019؛ Huber و همکاران 2019)
    \item ابزار: رگرسیون کوانتایل (Koenker \& Bassett 1978) — تابع زیانی با وزن $\tau$ برای
      کم‌زدن و $1-\tau$ برای زیادزدن
    \item کمینه‌کردن این زیان و حل مسئله‌ی Newsvendor یک کار واحدند، نه دو کار پشت سر هم
  \end{itemize}
  \begin{figure}
    \centering
    \includegraphics[width=0.68\textwidth]{slides_05_pinball}
  \end{figure}
  \note{مدل مستقیماً «عدد پخت» را می‌آموزد، نه «تقاضا» را. به همین دلیل مدل نهایی یک درخت
  گرادیان‌بوستینگ با هدف کوانتایلی است، همان انتخابی که در رقابت \lr{M5 Uncertainty} هم برنده
  شد. فرمول Pinball را بلند نگو، اسلاید پشتیبان دارد.}
\end{frame}

\begin{frame}{یک تصحیح که مال خود پروژه است}
  \begin{itemize}
    \item فرمول کلاسیک در فضای \textbf{تقاضا} تعریف شده؛ خروجی این پروژه چندکی از
      \textbf{نرخ عدم‌دریافت} است که متمم تقاضاست: $D = Res(1-\rho)$
    \item پس فرمول برمی‌گردد: $\displaystyle \tau^{*}_{\rho} = \frac{C_o}{C_u+C_o}$
    \item تأیید عددی با شبیه‌سازی مستقیم نیوزوندور روی ۴۰۰ هزار نمونه
  \end{itemize}
  \begin{figure}
    \centering
    \includegraphics[width=0.68\textwidth]{slides_06_tau_mapping}
  \end{figure}
  \note{این اسلاید را حذف نکن، تنها دستاورد نظری قابل‌دفاع پروژه است. سند اولیه‌ی مسئله این
  تبدیل را نداشت؛ با آن نسخه $\tau=0.10$ یعنی کمبود ارزان‌تر از مازاد، دقیقاً برعکس کل
  انگیزه‌ی پروژه. جمله‌ی پایانی: عدد $\tau$ خروجی مدل نیست، ورودی سیاستی است — مدیریت سلف
  نسبت $C_u/C_o$ را تعیین می‌کند و بقیه‌اش خودکار است. شاهد: ردیف ۳۶ decision\_log،
  \lr{reports/tau\_sensitivity.md}.}
\end{frame}

\begin{frame}{داده و قاعده‌ی برش}
  \begin{itemize}
    \item ۲ میلیون رزرو فردی، ۱۴۲ روز سرو، ۳۰ سلف، ۶ شهر
    \item تصمیم پخت ساعت ۱۵ روز قبل برای ناهار و ۲۳ برای شام گرفته می‌شود — مدل فقط چیزی
      را می‌بیند که آن لحظه واقعاً در دسترس بوده
  \end{itemize}
  \begin{figure}
    \centering
    \includegraphics[width=0.95\textwidth]{slides_07_timeline}
  \end{figure}
  \note{بگو: شام روز قبل، در لحظه‌ی تصمیمِ ناهار فردا هنوز سرو نشده، پس استفاده از آن نشتی
  است با اینکه تاریخش گذشته است. جدول کامل برش به پشتیبان می‌رود — این اسلاید را در ۲۰ ثانیه
  رد شو.}
\end{frame}
